Activation steering turns localized representations into control directions,
but localization alone does not reveal whether a direction has a selective
operating regime. We introduce \emph{Predictive Memory Localization} (PML),
which treats the measured-grid intervention path as the predictive object of
memory localization. PML separates random-calibrated target movement from
semantic-neighbor and capability damage, and compares static localization and
supervised geometry with a strength-disjoint low-dose causal response. Our
frozen study covers 3,000 records from nine datasets and fourteen domains,
yielding 30,000 distinct record--direction--layer paths and 210,000 distinct
path--strength evaluations. At layer 7, the geometry-derived RFM/AGOP direction reaches 13.1\% target-any and
12.3\% clean-any, exceeding random by 3.6 and 3.4 percentage points under a
record-paired bootstrap. Across record-, dataset-, and domain-grouped splits,
responses at $|\alpha|=0.1$ are the strongest signal for outcomes at disjoint strengths
$|\alpha| = 0.25$ or $0.5$. On held-out records, a predictor-driven selector
chooses a coefficient or abstains, improves utility and reduces
semantic-neighbor damage relative to a train-tuned fixed-strength policy, and
avoids most evaluations in a dense scan.
Across three residual-norm-matched base models, learned directions retain
selective-path gains and low-dose responses yield $0.801$--$0.828$
record-held-out macro AUROC.

PML therefore turns memory localization into a falsifiable forecast of
margin-level selective outcomes and a risk-aware intervention decision.
